\chapter{Marco normativo}
\label{ch:marco-normativo}

% Fuente canónica de títulos, fechas, vigencia y aplicación: Memoria/analisis/normativa-inventario.md
% (volcado literal de Bibliografia/Fuentes/Normativa/NormativaLinks.xlsx). Ninguna norma de este
% capítulo se cita con \cite ni figura en biblio.bib.
% TODO Mario: las referencias cruzadas a capítulos posteriores van en texto llano («el capítulo 3»)
% porque esos capítulos todavía no existen. Convertirlas en \ref{} conforme se vayan escribiendo.

La normativa decide antes que el cálculo. Qué configuraciones de autoconsumo con almacenamiento resultan admisibles, y cuáles quedan descartadas de partida, se resuelve en el marco regulatorio y no en la hoja de dimensionado, por lo que a continuación se relacionan las disposiciones aplicables al proyecto ordenadas por jerarquía: marco europeo, normas estatales con rango de ley, normas estatales de rango reglamentario e infrarreglamentario, normativa autonómica y local, y normas técnicas de referencia.

Cada disposición se identifica por su título completo en la primera mención y por su forma corta a partir de la segunda. De todas ellas se indica qué regulan y, sobre todo, de qué modo afectan a este proyecto en concreto, ya que el objeto del capítulo no consiste en resumir el ordenamiento sectorial sino en fijar las condiciones de contorno del diseño. Debido a dicho planteamiento la normativa no figura en el capítulo de referencias bibliográficas. Se relaciona íntegramente aquí.

Se evalúan tres variantes de implantación, y las tres reciben un tratamiento normativo distinto: una instalación sobre la cubierta del propio edificio, de 338 kWp; un conjunto de marquesinas sobre el aparcamiento situado a 216 m, de 680 kWp; y una planta en suelo de entre 2,00 y 2,50 MWp emplazada a 2,93 km del \ac{CPD}. Esa asimetría es el hilo conductor del capítulo. Separa lo que no pasa de instalación interior del consumidor de lo que exige permisos de acceso propios y pago por el uso de la red.

\section{Marco europeo}
\label{sec:norm-europeo}

Del marco europeo interesan aquí dos cosas, y ninguna de las dos impone obligaciones directas al titular del \ac{CPD}. Por un lado están los indicadores con los que se mide el rendimiento energético de un centro de proceso de datos, que son los que se emplean en el capítulo 3 para caracterizar el edificio. Por otro, la cobertura jurídica de los contratos de compra de electricidad a plazo, alternativa frente a la que se compara la variante de planta en suelo en el capítulo 9.

\subsection{Directiva (UE) 2023/1791, de 13 de septiembre, de eficiencia energética (refundición)}
\label{subsec:norm-directiva-2023-1791}

Refunde la normativa europea de eficiencia energética y somete por primera vez a los centros de proceso de datos a obligaciones específicas de seguimiento y publicación de su rendimiento energético e hídrico. Su artículo 12.1 exige a los Estados miembros que los titulares de centros con una potencia eléctrica demandada por los sistemas de tecnologías de la información igual o superior a 500 kW pongan a disposición del público la información de su anexo VII. La potencia media de tecnologías de la información del \ac{CPD} estudiado ronda los 295 kW. Queda por debajo del umbral y, por tanto, fuera de esa obligación.

Conviene precisar el alcance: se trata de una directiva, por lo que no resulta invocable frente a particulares mientras no se transponga al Derecho interno. La obligación alcanza al Estado, no directamente al titular de la instalación. A fecha de redacción de este trabajo la transposición del artículo 12 no se ha completado en España, de modo que la directiva se relaciona aquí como marco de referencia y como origen jurídico de los indicadores que se emplean más adelante, nunca como norma de cumplimiento directo.

Sin embargo, su artículo 11.2 sí resulta de interés para el titular. Impone auditoría energética periódica a los consumos medios anuales superiores a 10 TJ que no dispongan de sistema de gestión de la energía, y los 3.832 MWh anuales del centro equivalen a unos 13,8 TJ, por encima de dicho umbral. Cabe destacar además la letra d) de su anexo VI, que obliga a la auditoría a determinar el potencial de uso o producción rentable de energías renovables: exactamente el contenido de este trabajo.
% TODO Mario: confirmar si el umbral de 10 TJ del art. 11.2 se mide sobre el emplazamiento o sobre
% la empresa en su conjunto antes de afirmarlo en la defensa. Enlaza con la auditoría energética
% del CPD, que sigue sin confirmarse (memoria project_decisiones_bateria_auditoria).

\subsection{Reglamento Delegado (UE) 2024/1364, de 14 de marzo de 2024, relativo a la primera fase del establecimiento de un régimen de evaluación común de la Unión para centros de datos}
\label{subsec:norm-reglamento-2024-1364}

Desarrolla el artículo 33.3 de la directiva anterior y normaliza la medida de los indicadores de sostenibilidad de un centro de proceso de datos. Su anexo III define la eficacia en el uso de la energía (\ac{PUE}) como el cociente entre la energía total del centro y la energía de los sistemas de tecnologías de la información, y la eficacia en el uso del agua (\ac{WUE}) como el cociente entre el agua de entrada y esa misma energía. El anexo II precisa con qué norma técnica se mide cada término. No es un detalle menor, ya que de ello depende que los valores sean comparables entre instalaciones.

De este reglamento proceden las definiciones de \ac{PUE} y \ac{WUE} empleadas en el apartado de criticidad del capítulo 3, así como el criterio para computar la energía renovable propia. Dicho anexo II distingue la renovable acreditada mediante garantías de origen, la contratada mediante acuerdos de compraventa y la producida en el propio centro, por lo que la instalación en cubierta de este proyecto computaría en esta última categoría y un eventual contrato a plazo en la segunda. Por otro lado, contempla la declaración de las funciones de red que presta el almacenamiento, lo que ofrece un encaje normativo a la batería dimensionada en el capítulo 7.

\subsection{Reglamento (UE) 2024/1747, de 13 de junio, sobre la configuración del mercado de la electricidad de la Unión}
\label{subsec:norm-reglamento-2024-1747}

Reforma la configuración del mercado eléctrico europeo tras la crisis de precios de 2022 y da cobertura jurídica a los contratos de compra de electricidad a plazo (\acp{PPA}) y a los contratos bidireccionales por diferencias. Es un reglamento modificativo: su contenido se incorpora al articulado de los Reglamentos (UE) 2019/942 y (UE) 2019/943, por lo que al citar un precepto concreto procede referirse al texto consolidado de este último y no al propio 2024/1747.

Su relevancia para el trabajo es doble. Sustenta, en primer lugar, el \ac{PPA} como término de comparación económica frente al autoconsumo a distancia de la variante 3, comparación que se desarrolla en el capítulo 9. Obliga además a los gestores de red a publicar periódicamente la capacidad disponible para nuevas conexiones con detalle espacial, que es la base europea de los mapas de capacidad que deben consultarse para evaluar la conexión de la planta en suelo.

\section{Normativa estatal con rango de ley}
\label{sec:norm-ley}

Tres normas con rango de ley delimitan el marco español del autoconsumo y, en particular, las condiciones que hacen viable la variante de planta en suelo.

\subsection{Ley 24/2013, de 26 de diciembre, del Sector Eléctrico}
\label{subsec:norm-ley-24-2013}

Norma de cabecera del sector eléctrico español, de la que cuelga todo el desarrollo reglamentario que emplea este trabajo. Su artículo 9.1 define el autoconsumo y separa las modalidades sin excedentes, con mecanismo antivertido y un solo sujeto, y con excedentes, donde aparecen las figuras de consumidor y productor. El artículo 42 regula la línea directa, alternativa técnica evaluada para conectar las marquesinas del aparcamiento.

El precepto de mayor impacto económico es el artículo 9.5. Declara exenta de todo tipo de cargos y peajes la energía autoconsumida de origen renovable, cogeneración o residuos, pero añade a continuación que, cuando se produce transferencia de energía a través de la red de distribución en instalaciones próximas, pueden establecerse las cantidades que resulten de aplicación por el uso de dicha red. Ahí está la base legal del coste de red que soporta la variante 3, y que se cuantifica en el capítulo 9.

Desde el 22 de marzo de 2026 la ley incorpora dos reglas sobre almacenamiento que condicionan el diseño de la batería. El artículo 33.13 declara que los permisos de acceso y conexión de las instalaciones de almacenamiento son permisos de acceso flexibles desde la perspectiva de demanda. El artículo 39.3 establece, en su último párrafo, que no tienen consideración de consumidor las instalaciones de producción hibridadas con almacenamiento que, conforme a su permiso de acceso, puedan consumir energía de la red, con lo que se evita el doble pago sobre la misma energía. La hibridación propiamente dicha, en cambio, se regula en el artículo 33.12, que permite evacuar por el mismo punto de conexión y con la capacidad de acceso ya concedida.

\subsection{Real Decreto-ley 15/2018, de 5 de octubre, de medidas urgentes para la transición energética}
\label{subsec:norm-rdl-15-2018}

Norma fundacional del régimen actual de autoconsumo: suprimió el peaje de respaldo, consagró el autoconsumo colectivo y dio nueva redacción al artículo 9 de la Ley 24/2013. Para este trabajo importa sobre todo por su disposición adicional segunda, que declara exentas de obtener permisos de acceso y conexión para generación a las instalaciones de autoconsumo acogidas a la modalidad sin excedentes, así como a aquellas de potencia igual o inferior a 15 kW situadas en suelo urbanizado con las dotaciones y servicios requeridos por la legislación urbanística.

El matiz es decisivo y conviene explicitarlo. La exención de la primera letra opera por modalidad, no por potencia: una instalación sin excedentes está exenta cualquiera que sea su potencia, mientras que el umbral de 15 kW afecta únicamente a las modalidades con excedentes. Debido a ello las variantes de cubierta y de marquesinas quedan exentas pese a superar con holgura ese umbral. La cadena de citas correcta enlaza dicha disposición adicional segunda con el artículo 17 del Real Decreto 1183/2020, que constituye su desarrollo reglamentario y se remite a ella de forma expresa.

\subsection{Real Decreto-ley 7/2026, de 20 de marzo}
\label{subsec:norm-rdl-7-2026}

Norma de respuesta a la crisis de Oriente Medio que, en sus disposiciones finales de energía, reescribe las reglas de proximidad del autoconsumo. Es la que desbloquea la variante 3: sin ella, los 2,93 km entre la parcela y el \ac{CPD} incumplían el criterio general de 500 m por un factor próximo a seis. Fue convalidada por el Congreso el 26 de marzo de 2026 y se encuentra en vigor desde el 22 de marzo de ese mismo año.

Su disposición final decimocuarta actúa sobre el Real Decreto 244/2019 en dos apartados. El primero da nueva redacción al artículo 3.g).iii), mantiene los 500 m como regla general e incorpora un supuesto de 5.000 m para instalaciones fotovoltaicas o eólicas de hasta 5 MW conectadas por las redes de transporte o distribución. El segundo reescribe el artículo 4.5, cuya letra b) admite como única excepción la combinación de un autoconsumo individual sin excedentes con otro mediante instalaciones próximas a través de la red, y cuya letra c) obliga a que todo autoconsumo a través de la red pertenezca a la modalidad con excedentes.

Por su disposición final octava se introdujeron en la Ley 24/2013 el artículo 33.13 y el último párrafo del artículo 39.3 ya comentados, mientras que la decimoquinta reformó el Real Decreto 1183/2020 para crear la figura del acceso flexible. Su disposición adicional primera, en fin, habilita al Consejo de Ministros para aprobar por real decreto los requisitos de sostenibilidad de los centros de datos que se conecten a las redes, desarrollo que se comenta en el apartado~\ref{sec:intro-justificacion} por no constituir todavía derecho vigente.

Al manejar esta norma conviene advertir de dos extremos. El primero, que la variante 3 solo resulta viable con la redacción posterior al 22 de marzo de 2026, ya que las anteriores exigían cubierta de edificación o suelo industrial y una distancia máxima de 1.000 m, condiciones que una parcela rústica a 2.930 m no cumplía ni por emplazamiento ni por distancia. El segundo, que el Real Decreto-ley 7/2025 había introducido reglas semejantes pero quedó sin efecto al no ser convalidado. El contenido vigente procede exclusivamente de esta norma.

\section{Normativa estatal de rango reglamentario e infrarreglamentario}
\label{sec:norm-reglamento}

En este nivel se sitúan los reales decretos que desarrollan la Ley 24/2013 y, por debajo de ellos, las circulares y resoluciones de la Comisión Nacional de los Mercados y la Competencia (\ac{CNMC}), de naturaleza infrarreglamentaria al amparo del artículo 7 de la Ley 3/2013. Son estas últimas las que fijan la metodología y los valores de los peajes, así como las condiciones de acceso a las redes.

\subsection{Real Decreto 244/2019, de 5 de abril, de condiciones administrativas, técnicas y económicas del autoconsumo}
\label{subsec:norm-rd-244-2019}

Es la norma central del proyecto, ya que determina por sí sola la tipología jurídica de las tres variantes estudiadas. Su artículo 3.g) establece cuándo dos instalaciones se consideran próximas y distingue cuatro supuestos: conexión por la red interior del consumidor, unión mediante línea directa, situación dentro de la distancia máxima admisible e identidad de los catorce primeros dígitos de referencia catastral. El párrafo de cierre de dicho artículo precisa que el primer supuesto da lugar a instalaciones próximas de red interior, y los tres restantes a instalaciones próximas a través de la red.

La distancia se mide entre los equipos de medida en su proyección ortogonal en planta, no por el trazado de la línea de evacuación. Por otro lado, el artículo 3.h) define la potencia instalada de una instalación fotovoltaica como la potencia máxima del inversor, o la suma de las de los inversores, criterio determinante para acreditar que la planta en suelo se mantiene por debajo del límite de 5 MW del supuesto de 5.000 m.

El artículo 4 fija las modalidades de suministro con autoconsumo. Su apartado 2.a) condiciona la compensación simplificada de excedentes a que la potencia total de las instalaciones de producción asociadas no supere los 100 kW, umbral que la variante 3 rebasa con creces y que la deja fuera de dicho mecanismo. Los apartados 5.b) y 5.c), en la redacción de 2026, contienen la regla de combinación de modalidades y la obligación de excedentes ya descritas.

El régimen de permisos lo reparte el artículo 7.1.b). Su ordinal i) declara exentas de obtener permisos de acceso y conexión a las instalaciones sin excedentes; el ordinal iii) los exige, por cada instalación, a las instalaciones con excedentes de potencia superior a 15 kW. De ahí procede la asimetría que separa las variantes 1 y 2 de la variante 3, y que constituye la principal barrera de entrada de esta última, por encima de la superficie o del coste.

En el plano económico, los artículos 17.1 y 18.1 declaran la exención general de peajes y cargos sobre la energía autoconsumida. El artículo 17.5 la excepciona: cuando se produce transferencia de energía a través de la red de distribución en instalaciones próximas, los consumidores asociados deben satisfacer una cuantía por la utilización de dicha red, determinada por la \ac{CNMC}. El artículo 17.4.b) precisa, en cuanto al autoconsumo a través de la red, que el término de energía se factura sobre la energía horaria consumida de la red individualizada, criterio que se traslada al modelo horario del capítulo 7.

Por último, la disposición transitoria quinta resuelve el encaje de las baterías al exigir que los elementos de almacenamiento compartan equipo de medida y protecciones con la instalación de generación. En consecuencia la batería no constituye una instalación independiente ni solicita permiso propio, cuestión que se retoma en el capítulo 7.

La aplicación del real decreto al proyecto se resume así. La instalación en cubierta se acoge a la figura de instalación próxima de red interior y a la modalidad sin excedentes. Las marquesinas del aparcamiento se levantan sobre una referencia catastral que difiere de la del edificio a partir del séptimo dígito, por lo que no pueden ampararse en el supuesto de identidad catastral y se resuelven mediante línea propia al amparo del primer supuesto. La planta en suelo, en cambio, se configura como instalación próxima a través de la red y, en consecuencia, en modalidad con excedentes y con permisos de acceso y conexión propios.

\subsection{Real Decreto 1183/2020, de 29 de diciembre, de acceso y conexión a las redes de transporte y distribución}
\label{subsec:norm-rd-1183-2020}

Regula el procedimiento de acceso y conexión: quién solicita capacidad, en qué nudo, con qué plazos y con qué garantías. Es el itinerario administrativo de la variante 3 y confirma que las variantes 1 y 2 se lo ahorran, ya que su artículo 17 remite de forma expresa a la exención de la disposición adicional segunda del Real Decreto-ley 15/2018.

Su artículo 10.1 establece que, por encima de 100 kW, la solicitud se formula para un nudo o tramo de línea concreto de la red. El punto de conexión lo elige por tanto el promotor, no el gestor. Los apartados siguientes del mismo artículo fijan los plazos de subsanación y admisión que alimentan el cronograma de tramitación del capítulo 4, y el artículo 12 determina el contenido mínimo de la propuesta previa, incluidas las potencias de cortocircuito máxima y mínima de diseño, las condiciones de la línea de evacuación y el presupuesto económico pormenorizado.

El artículo 3.2 excluye del ámbito de aplicación a las instalaciones de almacenamiento cuando estas nunca inyectan energía a las redes de transporte o distribución. Debido a ello la batería situada tras el contador en las variantes 1 y 2, con antivertido, queda fuera del real decreto; otro tanto sucede con una batería hibridada con la planta en suelo que se cargue únicamente desde la instalación fotovoltaica. La frontera no la marca la ubicación física del equipo. La marca su capacidad de cargar desde la red.

Cabe destacar que este real decreto delimita el riesgo abierto más relevante del proyecto, que consiste en determinar si existe capacidad de acceso de generación disponible en un nudo razonablemente próximo a la parcela de la variante 3. Dicha comprobación debe resolverse antes de acometer la simulación detallada de esa variante.

\subsection{Real Decreto 148/2021, de 9 de marzo, sobre metodología de cálculo de los cargos del sistema eléctrico}
\label{subsec:norm-rd-148-2021}

Establece la metodología de cálculo de los cargos, componente de la factura que cubre los costes de política energética y que no debe confundirse con los peajes de red. Su artículo 1.2 aplica los cargos a los consumidores tal como los define el artículo 6.1.g) de la Ley 24/2013, mientras que el artículo 1.3.b) excluye del pago a la energía consumida por los titulares de instalaciones de almacenamiento que sea posteriormente inyectada en la red, con el fin de evitar el doble pago sobre la misma energía.

Su aplicación a este proyecto es acotada. Con la configuración adoptada, en la que la batería se carga exclusivamente desde la instalación fotovoltaica, no se produce retirada de energía de la red destinada al almacenamiento, por lo que la cuestión resulta indiferente. Conviene dejar constancia, sin embargo, de un extremo no resuelto: la exención se ancla a la condición de titular de instalación de almacenamiento del artículo 6.1.h) de la Ley 24/2013, y una batería situada tras el contador, cuyo titular es el \ac{CPD} en su condición de consumidor, no encaja de forma evidente en dicho supuesto.
% TODO Mario: cuestión abierta. Si el modelo económico llegase a contemplar carga desde red, hay
% que localizar criterio de la CNMC o declarar la exención como hipótesis explícita y llevarla al
% análisis de sensibilidad del capítulo 9.

\subsection{Real Decreto 1699/2011, de 18 de noviembre, de conexión a red de instalaciones de producción de pequeña potencia}
\label{subsec:norm-rd-1699-2011}

Es hoy una norma residual y muy vaciada, desplazada por el Real Decreto 244/2019 y por el Real Decreto 1183/2020. Conserva vigencia, no obstante, en un punto que afecta a este proyecto: su artículo 6, relativo al pliego de condiciones técnicas de conexión, permanece íntegramente vigente desde 2011, y el artículo 12.3.b) del Real Decreto 1183/2020 se remite a él para el pliego y el presupuesto de conexión de los autoconsumidores en su parte de generación.

Su procedimiento de acceso, en cambio, ya no rige. Los artículos 4.2 y 5, este último dedicado precisamente a la determinación de las condiciones técnicas de acceso y conexión, fueron derogados por el Real Decreto 1183/2020 con efecto de 31 de diciembre de 2020. Debido al alcance desigual de esas derogaciones es la norma de todo el capítulo con mayor riesgo de que se cite un precepto ya sin vigor, por lo que cada artículo se comprueba antes de llevarlo al texto.

\subsection{Circular 3/2020, de 15 de enero, de la CNMC, sobre metodología de peajes de transporte y distribución}
\label{subsec:norm-circular-3-2020}

Fija la metodología y la estructura de los peajes de acceso a las redes, y es la pieza que matiza la exención general del artículo 9.5 de la Ley 24/2013. Su artículo 2.3 declara aplicable la circular tanto a los pagos de los autoconsumidores por la energía consumida de la red como a los correspondientes a la energía autoconsumida en instalaciones próximas a través de la red. Son dos conceptos distintos y se facturan por separado.

El artículo 6.3 crea la familia de peajes TDA, integrada exclusivamente por términos de energía y sin término de potencia, y su letra c) aplica el peaje 6.1TDA a los suministros conectados en tensiones superiores a 1 kV e inferiores a 30 kV, correspondientes al nivel de tensión tarifario NT1. Dicha correspondencia entre el 6.1TDA y el NT1 es la que permite después leer los valores publicados, por lo que se explicita aquí. Por otro lado, el artículo 2.2.d) exceptúa del pago de peajes a las baterías de almacenamiento conectadas a la red de transporte o distribución, exención que alcanza a la batería como sujeto conectado a la red y no a la energía que el consumidor retira de ella.

La consecuencia para el proyecto es directa: la energía autoconsumida por la planta en suelo paga peaje de red, mientras que la de las variantes de cubierta y marquesinas está exenta. Los cargos, en cambio, siguen exentos en las tres variantes, ya que la salvedad legal se refiere únicamente al uso de la red de distribución. El artículo 13.1 habilita, en fin, la resolución anual que cuantifica los valores.

La circular se encuentra modificada por la Circular 1/2025, de 28 de enero, de la propia \ac{CNMC}, por lo que debe manejarse siempre su texto consolidado.
% TODO Mario: la disposición adicional primera de la Circular 3/2020 prevé revisión de la
% metodología cada seis años y el primer periodo regulatorio terminó el 31/12/2025. La CNMC declara
% la revisión «en fase final de tramitación». Reconfirmar el consolidado antes de cerrar el
% capítulo 9: es plausible que aparezca una circular sustitutiva antes de la entrega.

\subsection{Resolución de 18 de diciembre de 2025, de la CNMC, sobre los valores de los peajes de acceso aplicables desde el 1 de enero de 2026}
\label{subsec:norm-resolucion-peajes}

La circular anterior crea el peaje 6.1TDA pero no lo cuantifica. Los valores los fija esta resolución anual, dictada al amparo del artículo 13.1 de aquella y del artículo 7.1 bis de la Ley 3/2013, que es por tanto la fuente de los precios con los que se valora el coste de red de la variante 3, así como del peaje 6.1TD que el \ac{CPD} paga por la energía comprada a la red en las tres variantes. Ambos conjuntos de valores se recogen, con su desglose por periodos, en el capítulo 9.

Conviene advertir de una particularidad de lectura. La resolución no emplea en ningún punto la denominación 6.1TDA: indexa sus tablas por nivel de tensión tarifario, de NT0 a NT4. La correspondencia con el NT1 procede del artículo 6.3.c) de la Circular 3/2020 y debe explicitarse al citar la cifra, ya que de otro modo el valor parecería carecer de fuente. Por otro lado, la resolución correspondiente al ejercicio 2027 se publicará con posterioridad a la entrega de este trabajo, por lo que los valores de 2026 se consideran definitivos a efectos de cálculo.

\subsection{Circular 1/2024, de 27 de septiembre, de la CNMC, sobre acceso y conexión de las instalaciones de demanda}
\label{subsec:norm-circular-1-2024}

Regula por primera vez, mediante circular, el acceso y la conexión de la demanda y no solo de la generación. Su aportación central es la capacidad de acceso flexible, instrumento con el que una instalación puede conectarse en una red saturada renunciando a la garantía de suministro en todas las horas del año. El artículo 3.2 distingue únicamente dos tipos de capacidad de acceso: la firme u ordinaria, con garantía de suministro en todas las horas, y la flexible, en la que dicha garantía no existe.

Su aplicación a este proyecto es limitada, aunque ilustrativa. El \ac{CPD} se encuentra en explotación y dispone de contrato de \ac{ATR} vigente, por lo que no requiere una nueva solicitud como consumidor. Los párrafos finales del mismo artículo 3.2 impiden además solicitar u obtener capacidad de acceso flexible a los suministros que, conforme a la normativa aplicable, deban estar garantizados, lo que descarta esa vía para una carga crítica de operación continua. Sí resulta pertinente la figura, en cambio, para el almacenamiento, al que el artículo 33.13 de la Ley 24/2013 impone el acceso flexible por imperativo legal en el supuesto de que llegue a solicitar permiso propio.
% TODO Mario: no está verificado si la CNMC ha aprobado ya las especificaciones de detalle del
% art. 18 de esta circular, de las que depende que el acceso flexible sea operativo. Si el trabajo
% se apoya en la figura más allá de mencionarla, hay que resolverlo.

\subsection{Síntesis del régimen aplicable a cada variante}
\label{subsec:norm-sintesis}

Como se puede apreciar en la tabla~\ref{tab:norm-variantes}, las consecuencias normativas de las tres variantes divergen en todos los conceptos salvo en el de cargos. Son las que se trasladan al estudio de alternativas del capítulo 5 y al análisis económico del capítulo 9.

\begin{table}[H]
    \centering
    \caption{Consecuencias normativas de cada variante de implantación. Elaboración propia a partir del Real Decreto 244/2019, del Real Decreto-ley 15/2018 y de la Circular 3/2020 de la CNMC.}
    \label{tab:norm-variantes}
    \begin{tabular}{@{}p{2.9cm}p{3.0cm}p{3.0cm}p{3.6cm}@{}}
        \hline
        \textbf{Aspecto} & \textbf{V1. Cubierta} & \textbf{V2. Marquesinas} & \textbf{V3. Planta en suelo} \\
        \hline
        Figura del art. 3.g) & Próxima de red interior & Red interior mediante línea propia & Próxima a través de la red \\
        Modalidad del art. 4 & Sin excedentes & Sin excedentes & Con excedentes, obligatoria por el art. 4.5.c) \\
        Permiso de acceso y conexión & No exigible & No exigible & Exigible por cada instalación \\
        Peaje sobre la energía autoconsumida & Exenta & Exenta & Peaje 6.1TDA, nivel NT1 \\
        Cargos sobre la energía autoconsumida & Exenta & Exenta & Exenta \\
        Régimen de la batería & Comparte medida y protecciones con la generación & Comparte medida y protecciones con la generación & Igual, siempre que se cargue solo desde la instalación fotovoltaica \\
        \hline
    \end{tabular}
\end{table}

Procede, por último, una precisión de alcance. El proyecto de real decreto por el que se regulan los requisitos de sostenibilidad energética, medioambiental y de resiliencia y soberanía digital de los centros de datos, sometido a audiencia pública en agosto de 2026, se encuentra en tramitación y carece de valor jurídico. No se relaciona por ello entre la normativa aplicable a este proyecto, y su contenido se comenta en el apartado~\ref{sec:intro-justificacion} como argumento de actualidad del tema.

\section{Normativa autonómica y local}
\label{sec:norm-autonomica}

El nivel autonómico y local no afecta por igual a las tres variantes, y ni siquiera las reparte entre los mismos ayuntamientos. Las instalaciones sobre cubierta y sobre el aparcamiento se resuelven, en lo urbanístico, dentro de una parcela ya edificada y con uso consolidado del término de Tres Cantos, por lo que su tramitación se limita al título habilitante municipal. La planta en suelo es otra cosa: se implanta sobre una parcela rústica de Colmenar Viejo, ajena al recinto del \ac{CPD}, de modo que el régimen urbanístico del suelo determina si la variante resulta siquiera posible.

\subsection{Ley 9/2001, de 17 de julio, del Suelo, de la Comunidad de Madrid}
\label{subsec:norm-ley-9-2001}

Fija el régimen urbanístico del suelo en el ámbito autonómico. Su artículo 13 lo clasifica en urbano, urbanizable y no urbanizable de protección, y determina para cada clase qué usos resultan legítimos y bajo qué título habilitante. El artículo 26.1.c) admite de forma expresa la generación de energía destinada a autoconsumo como uso habilitable en suelo urbanizable no sectorizado, y el artículo 27.1.d) contempla en el mismo régimen los propios centros de proceso de datos.

Sin embargo, el artículo 16.1.b) califica como suelo no urbanizable de protección los terrenos que el planeamiento preserva únicamente por sus valores agrícolas, forestales o ganaderos. Para esa categoría el artículo 29.3 establece una lista tasada de usos autorizables en la que no figura la generación fotovoltaica en suelo, de modo que solo restaría la vía excepcional del artículo 29.1, condicionada a una legislación sectorial habilitante que el propio texto no llega a concretar.

Conviene acotar el alcance de esa restricción, ya que opera sobre el suelo no urbanizable de protección y no sobre toda parcela rústica. Que una finca figure catastrada con uso agrario no la clasifica: quien clasifica es el planeamiento municipal, y de él depende cuál de las dos vías resulta aplicable, la del artículo 26.1.c) o la del artículo 29.3. Debido a ello la clasificación urbanística de la parcela candidata pesa tanto como su superficie o su distancia al \ac{CPD} a la hora de seleccionarla, y así se trata en el capítulo 5.

El procedimiento aplicable en los supuestos de suelo no urbanizable de protección y de suelo urbanizable no sectorizado es la calificación urbanística regulada en los artículos 147 y 148 de la misma ley, con plazo de resolución de tres meses y sentido desestimatorio del silencio administrativo. Dicho plazo se incorpora al cronograma de tramitación del capítulo 4.

\subsection{Planeamiento municipal y ordenanzas}
\label{subsec:norm-planeamiento}

La parcela candidata de la variante 3 no está en Tres Cantos. Se consultaron su centroide y las cuatro esquinas de su rectángulo envolvente contra el servicio de consulta por coordenadas de la Sede Electrónica del Catastro, y los cinco puntos devuelven el mismo código de municipio, el 28045, correspondiente a Colmenar Viejo, con referencias catastrales rústicas de los polígonos 43 y 45. La misma consulta sobre la coordenada del \ac{CPD} devuelve Tres Cantos, lo que confirma que la diferencia entre ambos emplazamientos es real y no un artefacto del servicio. Las fincas afectadas son las parcelas 47, 213, 74 y 18 del polígono 45 y la parcela 188 del polígono 43.

El instrumento de planeamiento aplicable es, por tanto, la Modificación del Plan General de Colmenar Viejo, cuyo Título 5 regula el suelo no urbanizable y data de 1987. Es anterior a la Ley 9/2001 y emplea sus propias categorías: suelo no urbanizable Común, identificado con la letra N y desglosado en dos orientaciones de aprovechamiento, y suelo no urbanizable de Protección Especial, identificado con la letra P y desglosado en quince subáreas. Ningún plan posterior lo ha sustituido, ya que el municipio solo había alcanzado la fase de documento de avance en enero de 2025.

La traducción entre aquellas categorías y las de la ley vigente la resuelve la propia Ley 9/2001 en su disposición transitoria primera. Su letra c) ordena aplicar al suelo no urbanizable común el régimen que la ley establece para el suelo urbanizable no sectorizado; su letra d) remite el suelo no urbanizable especialmente protegido al régimen de suelo no urbanizable de protección. La disposición transitoria tercera añade, por su parte, que los planes aprobados definitivamente antes de su entrada en vigor mantienen su vigencia, por lo que dicha traducción opera aunque el municipio no haya completado la adaptación formal de su planeamiento.

De ahí resultan dos salidas, sin zona intermedia. Si la subparcela finalmente elegida está clasificada como N, la reconducción de la letra c) la sitúa en el régimen del suelo urbanizable no sectorizado y le resulta directamente aplicable el artículo 26.1.c), esto es, la generación de energía destinada a autoconsumo como uso habilitable mediante calificación urbanística, con el plazo de tres meses y el silencio desestimatorio ya señalados. Si está clasificada como P, la letra d) la lleva al artículo 29 y a la lista tasada de su apartado 3, donde la generación fotovoltaica en suelo no figura. La variante 3 es viable en el primer supuesto. Queda bloqueada en el segundo.

Cuál de los dos supuestos se da no puede afirmarse todavía. La distinción entre N y P la fijan los planos de zonificación a escala 1:10.000 del propio plan, que no forman parte del volumen de normas urbanísticas consultado y no se han podido localizar, por lo que se solicitan al Ayuntamiento de Colmenar Viejo. De ese único dato depende por completo la conclusión de este apartado. Queda sin localizar, por otro lado, la publicación de la aprobación definitiva de la modificación de 1987: lo consultado es la diligencia de ratificación municipal de 13 de febrero de 1987, con informe favorable de la Comunidad de Madrid. Treinta y nueve años de aplicación sin plan sustitutivo hacen razonable suponer que dicha aprobación se alcanzó, si bien esto puede deberse simplemente a la ausencia de un instrumento posterior y no equivale a una comprobación documental.

Sí se ha descartado, en cambio, la afección de espacio natural protegido. Los mismos cinco puntos se consultaron contra las capas oficiales de espacios naturales protegidos y de Red Natura 2000 del Inventario Español del Patrimonio Natural y de la Biodiversidad, y las cinco consultas devuelven resultado vacío en ambas capas. Como control se repitió la consulta sobre un punto de La Pedriza, que sí responde con el Parque Regional de la Cuenca Alta del Manzanares y con el lugar de la Cuenca del río Manzanares, de modo que las capas funcionan. La parcela no está dentro de ninguna de esas figuras. El resultado tiene, eso sí, un límite conocido: se comprobaron cinco puntos del rectángulo envolvente, no el contorno completo del polígono real.

Dos comprobaciones siguen abiertas, y procede dejar constancia de ellas. No se han verificado las vías pecuarias que pudieran cruzar la parcela, ya que no se localizó la capa cartográfica correspondiente de la Comunidad de Madrid; el texto del plan de 1987 relaciona cañadas, cordeles y veredas por nombre, sin que ninguno coincida de forma evidente con los parajes de las fincas candidatas, pero esa lectura no sustituye a una verificación geométrica. Tampoco se ha comprobado el criterio general 2 del anexo III de la Ley 21/2013, de 9 de diciembre, de evaluación ambiental, relativo al solape con elementos de infraestructura verde formalmente declarados. El criterio general 1 del mismo anexo, referido a espacios protegidos y a la Red Natura 2000, sí queda descartado por la verificación anterior.

En cuanto a las variantes 1 y 2, ambas se sitúan en término de Tres Cantos, cuyo Plan General de Ordenación Urbana fue aprobado en 2003. Al tratarse de instalaciones sobre una parcela ya edificada y con uso consolidado, la cuestión relevante no es la clasificación del suelo sino el título habilitante que exijan las ordenanzas municipales, extremo que se concreta en el capítulo 4.

% TODO Mario: pendiente único de este apartado, y el más crítico de todo el capítulo. Pedir al
% Ayuntamiento de Colmenar Viejo (Concejalía de Urbanismo) los planos de zonificación 1:10.000 del
% Plan General, para saber si las fincas candidatas («Arroyo del Buitre», «El Olivo», «La
% Vencosilla», «Los Llanillos») son suelo no urbanizable Común (N) o de Protección Especial (P).
% De ese dato depende que la variante 3 sea viable. Aprovechar la misma gestión para que confirmen
% si la Modificación de 1987 alcanzó aprobación definitiva y sigue siendo el instrumento vigente.
% Secundarios: vías pecuarias (capa cartográfica de la Comunidad de Madrid, no localizada);
% criterio general 2 del anexo III de la Ley 21/2013; y fecha y publicación en BOCM del Plan
% General de Tres Cantos de 2003. La Ley 21/2013 NO está en NormativaLinks.xlsx: darla de alta
% antes de desarrollarla más allá de la mención actual.

\section{Normas técnicas de referencia}
\label{sec:norm-tecnicas}

Junto a la normativa de autoconsumo y de acceso a la red, el diseño debe ajustarse a los reglamentos de seguridad industrial y a las normas técnicas de producto y de ensayo, que definen las condiciones de ejecución, las protecciones exigibles y las prestaciones garantizadas de los equipos. Dichas normas no alteran la modalidad de autoconsumo ni el régimen económico. Sí determinan, en cambio, el dimensionado eléctrico, las distancias de seguridad, las características de la aparamenta y los criterios de aceptación de módulos, inversores y baterías.

En este trabajo se emplean en tres puntos concretos: el dimensionado eléctrico y el esquema unifilar del capítulo 6, la selección de componentes del mismo capítulo y la verificación de las garantías de rendimiento del capítulo 8. Su relación se cierra una vez fijada la solución adoptada, ya que el conjunto aplicable depende del nivel de tensión del punto de conexión y de la tipología constructiva finalmente seleccionada.

Dos normas quedan ya fijadas por remisión del Reglamento Delegado (UE) 2024/1364. La medida de la energía total del centro y de la energía de los sistemas de tecnologías de la información se realiza conforme a la norma EN 50600-4-2, y la del agua de entrada conforme a la norma EN 50600-4-9. Ambas pertenecen a la serie europea de infraestructura de centros de proceso de datos y son las que hacen comparable el valor de \ac{PUE} del edificio con los umbrales de referencia manejados en el capítulo 3.

% TODO Mario: resto del apartado pendiente. Otro agente investiga en paralelo qué normas
% corresponden. Candidatas que hay que verificar y dar de alta en NormativaLinks.xlsx ANTES de
% redactarlas (no escribirlas hasta tener título completo, edición y vigencia):
%   - Reglamento electrotécnico para baja tensión y sus instrucciones técnicas complementarias.
%   - Reglamento sobre condiciones técnicas y garantías de seguridad en instalaciones eléctricas de
%     alta tensión y sus ITC, aplicable al centro de transformación de la variante 3.
%   - Código Técnico de la Edificación, documentos básicos DB-SE (acciones en la edificación: viento
%     y nieve sobre cubierta y marquesinas) y DB-HE (ahorro de energía).
%   - Reglamento unificado de puntos de medida del sistema eléctrico, para la configuración de
%     medida de cada variante.
%   - Normas UNE-EN e IEC de módulos fotovoltaicos (cualificación de diseño, seguridad, medida de
%     potencia) y de inversores conectados a red.
%   - Resto de la serie EN 50600 aplicable a redundancia y criticidad (capítulo 3).
%   - Normativa de baterías y de residuos de aparatos eléctricos y electrónicos, para el apartado de
%     fin de vida del capítulo 10.
% Decidir con el tutor si los reglamentos de seguridad industrial se mantienen aquí o se trasladan
% al apartado 2.3, ya que formalmente tienen rango reglamentario.
